% !TEX root = ../notes.tex

\documentclass[../notes.tex]{subfiles}

\begin{document}

\section{October 1}
This talk was given by Jeremey Hahn. We will discuss quasicoherent sheaves.

\subsection{Stefanich Rings}
As usual, we denote by $\op{Pr}^L_\kappa$ the category of $\kappa$-presentable categories, $\kappa$ is some regular cardinal. We will mostly take $\kappa\coloneqq\aleph_1$. The morphisms are given by left adjoints which preserve $\kappa$-compact objects.
\begin{remark}
	Professor Hahn's $\aleph$s are immaculate.
\end{remark}
\begin{remark}
	It is convenient to use $\aleph_1$ instead of (say) $\aleph_0$ because $\aleph_1$-filtered colimits commute with countable limits, and countable limits appear frequently. For example, one may want to take the totalization of a cosimplicial object, which is a countable limit.
\end{remark}
\begin{example}
	Let $\mc C$ be the homotopy category of derived $p$-complete abelian groups. (In short, $\mc C$ is the category of abelian groups which are isomorphic to their own $p$-completions.) For example, $\ZZ_p$ is a projective generator of $\mc C$. However, $\ZZ_p$ is not $\omega$-compact: taking $\op{Hom}(\ZZ_p,-)$ on the diagram
	\[\colim\left(\ZZ_p\stackrel p\to\ZZ_p\stackrel p\to\ZZ_p\stackrel p\to\cdots\right).\]
	On one hand, $\op{Hom}(\ZZ_p,\colim\ZZ_p)=\op{Hom}(\ZZ_p,\QQ_p^\land)=0$. On the other hand, $\colim\op{Hom}(\ZZ_p,\ZZ_p)=\QQ_p$.
\end{example}
Recall that $\op{Pr}^L_\kappa$ admits a tensor product. We will be interested in commutative algebra objects $\mc C$ in $\op{Pr}^L_\kappa$, which admit module categories $\op{Mod}_{\mc C}\op{Pr}^L_\kappa$.
\begin{remark}
	Roughly speaking, a category in $\op{Mod}_{\mc C}\op{Pr}^L_\kappa$ is enriched over $\mc C$.
\end{remark}
In general, module categories acquire a symmetric monoidal structure. In other words, $\mathrm{Mod}_{\mc C}\op{Pr}^L_\kappa$ is a commutative algebra object in $\op{Pr}^L_\kappa$.
\begin{theorem} \label{thm:upgrade-pr}
	There is a fully faithful map $\op{CAlg}(\mc C)\to\op{CAlg}(\op{Mod}_{\mc C}\op{Pr}^L_\kappa)$ (in $\op{Pr}^L_\kappa$) given by sending $A$ to $\op{Mod}_A\mc C$.
\end{theorem}
\begin{remark} \label{rem:upgrade-pr}
	To be a morphism, this map needs a right adjoint. The right adjoint sends $\mc D$ to $\op{Hom}_{\mc D}(1)$. Because $\mc D$ is $\mc C$-linear, $\op{Hom}_D(1)$ upgrades to a $\mc C$-linear commutative algebra object.
\end{remark}
We are now ready to define a tower.
\begin{definition}
	Define $1\op{Pr}\coloneqq\op{Pr}^L_{\aleph_1}$. Then for each $n\ge1$, define $(n+1)\op{Pr}\coloneqq\op{Mod}_{n\mathrm{Pr}}(1\mathrm{Pr})$. Intuitively, these are the $n\mathrm{Pr}$-linear presentable categories.
\end{definition}
% \begin{remark}
% 	Here, we are applying \Cref{thm:upgrade-pr} to define these higher $\op{Pr}$s. For example, $2\op{Pr}$ is a commutative algebra object in $1\op{Pr}$, so $\op{Mod}_{2\op{Pr}}(1{\op{Pr}})$ is a commutative algebra object in $\op{Mod}_{1\mathrm{Pr}}(1\mathrm{Pr})$
% \end{remark}
\begin{remark}
	There are fully faithful functors
	\[\op{CAlg}(1\mathrm{Pr})\to\op{CAlg}(2\mathrm{Pr})\to\cdots\]
	provided by \Cref{thm:upgrade-pr}.
\end{remark}
This lets us cleanly define Stefanich rings.
\begin{definition}
	Define the category of \textit{Stefanich rings} to be the $\aleph_1$-presentable symmetric monoidal category
	\[\colim\left(\op{CAlg}(1\mathrm{Pr})\to\op{CAlg}(2\mathrm{Pr})\to\cdots\right),\]
	where the colimit takes place in $1\mathrm{Pr}$.
\end{definition}
\begin{remark}
	Colimits in $1\mathrm{Pr}$ have underlying $\infty$-category given by the limit in just the category of $\infty$-categories (of the right adjoints). It follows that Stefanich rings can be realized as sequences $(A_1,A_2,\ldots)$ where $A_n\in\op{CAlg}(n\mathrm{Pr})$, and we have chosen isomorphisms $A_{n-1}\cong\op{End}_{A_n}(1)$. (This latter isomorphism upgrades \Cref{rem:upgrade-pr}.)
\end{remark}

\subsection{Fun with \texorpdfstring{$\mathbb E_\infty$}{ Einf}-Rings}
Here is an $A$-linear variant.
\begin{example}
	Suppose $A$ is an $\mathbb E_\infty$-ring. Then $\mathrm{Mod}_A(\mathrm{Spaces})=D(A)$, which is a commutativ algebra object in $1\mathrm{Pr}$. Accordingly, we define $1\mathrm{Pr}_A\coloneqq\mathrm{Mod}_{D(A)}(1\mathrm{Pr})$ and then continue inductively by $(n+1)\mathrm{Pr}_A\coloneqq\mathrm{Mod}_{n\mathrm{Pr}}(1\mathrm{Pr})$.
\end{example}
\begin{remark} \label{rem:a-is-stefanich}
	Note that $n\mathrm{Pr}_A=\mathrm{Mod}_{(n-1)\mathrm{Pr}_A}(1\mathrm{Pr})$. There is a natural map of commutative algebra objects $(n-1)\mathrm{Pr}\to(n-1)\mathrm{Pr}_A$, so we can nest our modules to make this category
	\[\mathrm{Mod}_{(n-1)\mathrm{Pr}_A}(\mathrm{Mor}_{(n-1)\mathrm{Pr}}(1\mathrm{Pr}))=\mathrm{Mod}_{(n-1)\mathrm{Pr}_A}(n\mathrm{Pr}).\]
	This is now a commutative algebra object in $n\mathrm{Pr}$. Thus, $\{n\mathrm{Pr}_A\}_n$ is a Stefanich ring.
\end{remark}
\begin{remark}
	For each $n\ge0$, the map sending an $\mathbb E_\infty$-ring $A$ to $n\mathrm{Pr}_A\in\op{CAlg}(1\mathrm{Pr})$ upgrades to a functor.
\end{remark}
\begin{remark}
	It turns out that $n\mathrm{Pr}_A$ has a $\otimes$-$\op{Hom}$ adjunction. Also, for a map $f\colon A\to B$ of $\mathbb E_\infty$-rings, there is a pullback functor $f_n^*\colon n\mathrm{Pr}_A\to n\mathrm{Pr}_B$. (This is basically a tensor product. We already met this map in \Cref{rem:a-is-stefanich}.) There is also a functor $f_{n*}$, which is basically restriction.
\end{remark}
Here are some more difficult elaborations on this remark. All of these results are proved by induction on $n$. Typically, the base case is harder than the induction.
\begin{proposition}
	Let $f\colon A\to B$ be a map of $\mathbb E_\infty$-rings. Then $f_n^*\colon n\mathrm{Pr}_A\to n\mathrm{Pr}_B$ admits an $n\mathrm{Pr}_A$-linear colimit-preserving right adjoint $f_{n*}$. Moreover, $f_{n*}$ commutes with base-change and satisfies a projection formula.
\end{proposition}
\begin{proof}
	Let's argue this for $n=0$. Then we are on the hunt for a functor $f_*\colon D(B)\to D(A)$, which we construct as restriction of scalars along $f$. This provides the adjoint, so it only remains to check that $f_*$ preserves colimits; this holds because $B$-modules can already be computed in $\mathrm{Spectra}$.
\end{proof}
\begin{definition}
	A category $\mc C$ is \textit{countably presented} if and only if it is $\aleph_1$-compactly generated.
\end{definition}
\begin{proposition} \label{prop:get-shriek}
	If $B$ is a countably presented $A$-algebra, then $f_{n*}$ preserves $\aleph_1$-compact objects.
\end{proposition}
\begin{proof}
	Let's start with $n=0$. Let $D(B)^{\kappa}$ denote the set of compact objects. This turns out to be the subcategory generated by $B$ under countable colimits and shifts. Thus, $f_{0*}$ preserving compact objects is equivalent to checking that $B$ is $\kappa$-compact as an $A$-module. Well, because $B$ is countably presented, we can realize $B$ as a countable colimit of algebras of the form
	\[\op{Free}_A(M)=\bigoplus_{n\in\NN}M^{\otimes_An}_{h\Sigma_n},\]
	where $M$ is a compact $A$-module.

	In order to have done this at least once, we indicate how to do the induction on $n$. By the induction, it turns out that
	% It basically amounts to show that $n\mathrm{Pr}_A\to n\mathrm{Pr}_B$ makes $n\mathrm{Pr}_B$ countably presented as an $n\mathrm{Pr}_A$-module. (The case $n=0$ is the hypothesis, and the case $n=1$ was handled in the previous paragraph.) Thus, for our induction,
	we may suppose $(n-1)\mathrm{Pr}_A\to(n-1)\mathrm{Pr}_B$ is countably presented. Now, we must show that the map $f_{n*}\colon n\mathrm{Pr}_B\to n\mathrm{Pr}_A$ sends compact objects to compact objects. It turns out that the compact objects are generated by objects of the form $X\otimes(n-1)\mathrm{Pr}_B$ where $X$ is compact in $n\mathrm{Pr}$, so it basically suffices to handle $n\mathrm{Pr}_B$. Then we apply an argument similar to the one of the previous paragraph.
\end{proof}
\begin{remark}
	In the setting of \Cref{prop:get-shriek}, we may write $f_n^!$ for the additional right adjoint of $f_{n*}$.
\end{remark}
\begin{lemma} \label{lem:pr-commutes-limits}
	The functor $A\mapsto n\mathrm{Pr}_A$ commutes with finite limits.
\end{lemma}
\begin{proof}
	At $n=0$, this is basically asking us to show that $D(A\times B)=D(A)\times D(B)$. This is true because $A\times B$ admits idempotents which allows us to recover any module as a product of an $A$-module and a $B$-module.

	We will not do the whole induction, but let's indicate $n=1$. This amounts to checking
	\[\mathrm{Mod}_{R\times S}(1\mathrm{Pr})=\mathrm{Mod}_R(1\mathrm{Pr})\times\mathrm{Mod}_S(1\mathrm{Pr}).\]
	Thus, we are trying to say that a category which is $(R\times S)$-linear admits linearity by $R$ and $S$ separately. The moral is that $R$ and $S$ themselves are idempotent objects, so we can take tensor products to exhibit these equivalences.
\end{proof}
\begin{lemma} \label{lem:dual-of-pr}
	Let $f\colon A\to B$ be a countably presented map of $\mathbb E_\infty$-algebras. For $n\ge1$, the object $f_{n*}(1)\in n\mathrm{Pr}_A$ is dualizable.
\end{lemma}
\begin{proof}
	This is an application of Morita theory. Note $f_{n*}(1)=(n-1)\mathrm{Pr}_B$. In fact, this object is self-dual. The evaluation map needs to map out of $(n-1)\mathrm{Pr}_B\otimes(n-1)\mathrm{Pr}_B$ (where the tensor product takes place in $n\mathrm{Pr}_A$). Taking modules preserves colimits, so this is just $(n-1)\op{Pr}_{B\otimes_AB}$; there is then a diagonal map to $(n-1)\mathrm{Pr}_B$, which we then push forward to $(n-1)\mathrm{Pr}_A$.

	Analogously, the co-evaluation functor is given by
	\[(n-1)\mathrm{Pr}_A\stackrel{f^*}\to(n-1)\mathrm{Pr}_B\stackrel{\Delta_*}\to(n-1)\mathrm{Pr}_{B\otimes_AB}\cong(n-1)\mathrm{Pr}_B\otimes_{(n-1)\mathrm{Pr}_A}(n-1)\mathrm{Pr}_B.\]
	One then checks that these evaluations and co-evaluations make $(n-1)\mathrm{Pr}_B$ dualizable.
\end{proof}
\begin{remark}
	This may be false for $n=0$: indeed, we are asking for a countably presented ring to be a finitely presented modules. (A dualizable module is a finitely presented module.)
\end{remark}

\subsection{Descent}
We are now ready to prove something a little harder.
\begin{theorem}[Descent] \label{thm:descent}
	For each $n\ge0$, the functor from $\mathbb E_\infty$-rings to $1\mathrm{Pr}$ defined by $A\mapsto n\mathrm{Pr}_A$ is a hypersheaf for the fpqc topology.
\end{theorem}
\begin{remark}
	One could also try to get a hypersheaf to $\op{CAlg}((n-1)\mathrm{Pr})$. This turns out to be the same because the forgetful functor preserves limits.
\end{remark}
We should probably define some of these words. Let's start with fpqc.
\begin{definition}[faithfully flat]
	A map $f\colon A\to B$ of $\mathbb E_\infty$-rings is \textit{faithfully flat} if and only if
	\begin{itemize}
		\item $\pi_0f$ is faithfully flat, and
		\item the induced map $\pi_iA\otimes_{\pi_0A}\pi_0B\cong\pi_iB$ is an isomorphism for all $i\in\ZZ$.
	\end{itemize}
\end{definition}
The following explains why we will take fpqc to basically mean faithfully flat.
\begin{lemma} \label{lem:reduce-to-countably-presented}
	Let $f\colon A\to B$ be a faithfully flat map of $\mathbb E_\infty$-rings. Then $f$ is an $\aleph_1$-filtered colimit of countable presented faithfully flat maps $A\to B_i$.
\end{lemma}
\begin{proof}
	Let's wave our hands. For classical rings, this is due to Waterhouse. One starts by doing this for modules, and then one uses the module construction to build some algebras.
\end{proof}
Let's go ahead and engage with \Cref{thm:descent}.
\begin{proof}[Proof of \Cref{thm:descent}]
	By \Cref{lem:pr-commutes-limits}, $A\mapsto n\mathrm{Pr}_A$ commutes with finite products. Thus, it is enough to check that any fpqc hypercover $A\to A_\bullet$ induces an isomorphism
	\[n\mathrm{Pr}_A\to\lim_\Delta n\mathrm{Pr}_{A_\bullet}\]
	in $1\mathrm{Pr}$. By an elaboration of \Cref{lem:reduce-to-countably-presented}, every hypercover is an $\aleph_1$-filtered colimit of hypercovers which are level-wise countably presented $\mathbb E_\infty$-$A$-algebras. Thus, we may assume that each $A_\bullet$ in our cosimplicial object is a countably presented $A$-algebra.

	As usual, we will induct. Let's start with $n=0$. We want to show that the natural map
	\[D(A)\to\lim_\Delta D(A_\bullet)\]
	is an equivalence. The internal maps in the limit are given by $f^*$s. Because each $A_\bullet$ is countably presented, $D(A_\bullet)$ is self-dual by \Cref{lem:dual-of-pr}, so it is enough to show the dual of this map is an isomorphism. Namely, we want to show that the natural map
	\[\colim_\Delta D(A_\bullet)\to D(A)\]
	is an isomorphism, where the colimit is taken over $f_*$s. Now, $f^!$ is the right adjoint of $f_*$, so it is enough to show that the natural map
	\[D(A)\to\colim_\Delta D(A_\bullet)\]
	is an isomorphism, where the limit is taken over $f^!$s. Namely, $\colim_\Delta^*D(A_\bullet)=\lim^!_\Delta D(A_\bullet)$.
	
	One now shows that this functor is fully faithful and essentially surjective separately. We will only discuss that it is fully faithful. Well, it admits a left adjoint given by the geometric realization of restriction of scalars along $D(A_\bullet)\to D(A)$. By a unit argument, it is equivalent to checking that the induced map
	\[\colim_{\Delta\opp}\op{Hom}_A(A_\bullet,M)\to M\]
	is an isomorphism. By taking $\op{Hom}(-,M)$, it is enough to show that $A$ is isomorphic to $\{\op{Tot}^mA_\bullet\}_m$ as pro-objects, which turns out to follow from the hypersheaf condition: the cofiber of $A\to\op{Tot}^mA_\bullet$ is ghost, and the composite of two ghost maps is null.

	For the induction, we require the following.
	\begin{theorem}[Ambidexterity] \label{thm:ambi}
		Let $f\colon A\to B$ be a countably presented map of $\mathbb E_\infty$-algebras. For $n\ge1$, the functor $f_n^*\colon n\mathrm{Pr}_A\to n\mathrm{Pr}_B$ admits an $n\mathrm{Pr}_A$-linear left adjoint, which is isomorphic to $f_{n*}$.
	\end{theorem}
	The induction is now formal. Due to time, we omit details.
\end{proof}
Let's discuss how one can produce the ambidexterity of \Cref{thm:ambi}.
\begin{proposition} \label{prop:get-ambi}
	Let $\CC$ be an $(\infty,3)$-category, and let $F\colon X\to Y$ be a morphism in $\CC$. Then the following are equivalent:
	\begin{listroman}
		\item $F$ admits a right adjoint $F^R\colon Y\to X$, and ${\id_X}\to F^RF$ admits a right adjoint in $\op{End}_\CC(X)$, and $FF^R\to\id_Y$ admits a right adjoint in $\op{End}_\CC(Y)$; and
		\item $F$ admits a left adjoint $F^L\colon Y\to X$, and ${\id_X}\to F^RF$ admits a left adjoint in $\op{End}_\CC(X)$, and $FF^R\to\id_Y$ admits a left adjoint in $\op{End}_\CC(Y)$.
	\end{listroman}
	In this case, $F^L\cong F^R$.
\end{proposition}
\begin{proof}
	Formal nonsense. The idea is that units becomes counits. According to Professor Hahn, this is a Five lemma for $(\infty,3)$-categories.
\end{proof}
\begin{proof}[Proof of \Cref{thm:ambi}]
	We use \Cref{prop:get-ambi}. For $n\ge1$, $\CC\coloneqq(n+1)\mathrm{Pr}_A$ is an $(\infty,3)$-category. Then we will check (a) for the functor $f_n^*\colon n\mathrm{Pr}_A\to n\mathrm{Pr}_A$. We already know this has a right adjoint, and then we need to exhibit some more right adjoints, which is the purpose of countability hypotheses.
\end{proof}

\end{document}